\documentclass[11pt]{article}
\usepackage[margin=25mm]{geometry}
\usepackage{fontspec}\setmainfont{Latin Modern Roman}
\usepackage{amsmath,amssymb,mathtools}
\usepackage{unicode-math}\setmathfont{Latin Modern Math}
\usepackage{xurl,hyperref}
\hypersetup{colorlinks=true,urlcolor=blue}
\setcounter{secnumdepth}{0}\setlength{\parindent}{0pt}\setlength{\parskip}{5pt}\setlength{\emergencystretch}{3em}
\newcommand{\sref}[2]{\hyperlink{source:#1:#2}{[S#2]}}
\newcommand{\cataloguescope}[1]{\paragraph{Recorded target.}#1}
\begin{document}
\section{Kourovka Notebook Problem 21.106}\label{statement}

\subsection{Definitions and mathematical statement}

\paragraph{Shared notation.}$\mathbb N=\{1,2,\ldots\}$, and $\mathbb Z,\mathbb Q,\mathbb R,\mathbb C$ denote the integers, rationals, reals and complex numbers. Any other conventions are specified locally.


The group language has multiplication, inverse, and the identity constant. A first-order formula $\phi(x)$ has one free variable and no parameters from the group; other variables may be bound by quantifiers. Put $G_\phi=\{g\in G:G\models\phi(g)\}$ and $\phi(G)=\langle G_\phi\rangle$. A group $G$ is residually finite if for each $g\ne1$ some homomorphism from $G$ to a finite group has nonidentity image of $g$. A formula is concise on a class of groups when finiteness of $G_\phi$ always implies finiteness of $\phi(G)$.
\paragraph{Statement recorded in the supplied source.}
Is every parameter-free first-order group formula with one free variable concise on the class of all residually finite groups? In symbols, must
\[
 \bigl(G\text{ residually finite and }|G_\phi|<\infty\bigr)
 \quad\Longrightarrow\quad|\langle G_\phi\rangle|<\infty
\]
hold for every such $\phi$ and every such $G$, with no finite-generation or torsion assumption? \sref{statement}{2}
\subsection{Short English statement}
Do finitely many values of any one-variable first-order formula always generate a finite subgroup inside a residually finite group?
\subsection{Sources}
\begin{description}
\item[\hypertarget{source:statement:1}{S1}] ULAM.AI and the UnsolvedMath Contributors, \emph{UnsolvedMath: A Curated Collection of Open Mathematics Problems}, catalogue entry. CC BY 4.0. \url{https://huggingface.co/datasets/ulamai/UnsolvedMath}.
\item[\hypertarget{source:statement:2}{S2}] M. Petschick, Problem 21.106, in E. I. Khukhro and V. D. Mazurov (eds.), \emph{The Kourovka Notebook: Unsolved Problems in Group Theory}, 21st edition, arXiv:1401.0300v48 (6 October 2026). This version records a negative answer; the displayed question is the historically registered target. \url{https://arxiv.org/abs/1401.0300v48}.
\end{description}
\label{end:statement}
\end{document}
